\chapter{Invariant, Presentation, and Event Identity}
\label{ch:essay-presentation-identity}

The preceding chapters have treated a projection as a map from a latent structure to an observable realized in some physical channel, and have asked when that projection preserves admissibility structure. This chapter adds a distinction that the projection question, posed at that level of generality, does not yet make room for: the difference between two realizations of the same structure being \emph{another instance} of it and being \emph{a reference back to} an instance that already occurred. A recording of an utterance replayed a thousand times and a thousand independent speakers uttering the same sentence produce, under the apparatus developed so far, exactly the same observable if the observable is defined by structure and channel alone. Whether a memory or evidential system ought to treat these two cases identically is not a question the earlier chapters answer, because the apparatus as stated has no term for the distinction. This chapter supplies one, and shows that the distinction is not an addition external to the framework but a further instance of the observation map machinery already proven in Chapter~4 of \emph{Admissible Reconstruction}.

\section{Three Notions of Sameness}

\begin{definition}[Invariant, presentation, and occurrence]
\label{def:presentation-identity}
Let $O$ denote an admissibility-equivalence class in the sense of Chapter~1 of \emph{Admissible Reconstruction}: the class of all representations related by some admissibility-isomorphism $\varphi$, as established by the remark following Theorem~1.1 that representation invariance is itself an equivalence relation. Let $\theta$ denote a presentation parameter --- voice, font, accent, articulatory channel, or any other substrate-specific datum excluded from semantic content by Corollary~1.2. Write $P = \mathrm{realize}(O, \theta)$ for a particular presentation of $O$ under parameters $\theta$: a specific spoken, written, or gestured token. Let $\varepsilon$ denote an \emph{occurrence}: a location in physical time and history at which some presentation was actually produced, together with a provenance pointer that is either \textsc{original}, meaning the occurrence is itself the production of $P$, or a reference to an earlier occurrence $\varepsilon'$, meaning the occurrence is a replay, quotation, or retrieval of $P$ as originally produced at $\varepsilon'$.
\end{definition}

Three questions can now be distinguished that Definition~\ref{def:presentation-identity} keeps formally separate rather than allowing one to stand in for the other three. \emph{Invariant identity} asks whether two tokens realize the same $O$: whether ``give me the book'' and ``hand me the book'' express the same request, which Chapter~1's apparatus answers by exhibiting an admissibility-isomorphism between them, as Chapter~\ref{ch:essay-worked-isomorphism} did explicitly for a finite reduction. \emph{Presentation identity} asks the finer question of whether two tokens realize the same $O$ under the same $\theta$: the same request, in the same voice, at the same rate, in the same font. \emph{Event identity} asks whether two occurrences $\varepsilon_1$ and $\varepsilon_2$ are the same occurrence, or one is a reference to the other, or they are independent originals that happen to realize the same $(O, \theta)$. These three questions are logically independent: two tokens can share $O$ while differing in $\theta$ and in $\varepsilon$; can share $O$ and $\theta$ while differing in $\varepsilon$, as when the same recording is replayed at two different times; or, in the case this chapter is concerned with, can share $O$ and $\theta$ through entirely distinct and unrelated occurrences, as when two speakers, unknown to each other, happen to produce the identical utterance in the identical voice-independent structure.

\section{The Collapse as an Observation Map}

A memory or evidential system does not have direct access to $\varepsilon$; it has access only to what it recorded at the time, and what it recorded is itself the output of some observation map. Consider a system whose recording procedure is $\mathrm{obs}(\varepsilon) = (O(\varepsilon), \theta(\varepsilon))$: it stores the invariant and the presentation parameters of whatever it observes, and discards the occurrence's provenance pointer, retaining no record of whether a given presentation is an original production or a reference to an earlier one. Such a system cannot subsequently tell, from its own records, whether a given $(O, \theta)$ pair was produced by ten thousand independent occurrences or is the same occurrence recorded, retrieved, and re-presented ten thousand times.

\begin{proposition}
\label{prop:presentation-collision}
Let $H$ be a complete history of occurrences, each tagged with $(O, \theta, \varepsilon)$, and let $\mathrm{obs}$ be any observation map of the form $\mathrm{obs}(H) = \psi\bigl(\{(O(\varepsilon), \theta(\varepsilon)) : \varepsilon \in H\}\bigr)$ for some function $\psi$ --- that is, any observation map that records the multiset of invariant-presentation pairs occurring in $H$ but not the provenance pointers distinguishing originals from references. Then there exist histories $H_1 \neq H_2$, differing only in whether a given $(O,\theta)$ pair arose from one occurrence referenced many times or from many independent original occurrences, such that $\mathrm{obs}(H_1) = \mathrm{obs}(H_2)$.
\end{proposition}

\begin{proof}
Let $H_1$ consist of a single original occurrence $\varepsilon_0$ realizing $(O, \theta)$, together with $n - 1$ further occurrences each carrying a provenance pointer to $\varepsilon_0$. Let $H_2$ consist of $n$ occurrences, each realizing $(O, \theta)$ and each tagged \textsc{original}, with no provenance pointer relating any of them. By construction $H_1 \neq H_2$, since the occurrences differ in provenance structure even though every occurrence in both histories realizes the same $(O, \theta)$. Since $\mathrm{obs}$ is defined as a function of the multiset $\{(O(\varepsilon), \theta(\varepsilon)) : \varepsilon \in H\}$ alone, and this multiset is identical for $H_1$ and $H_2$ --- both consist of $n$ copies of the pair $(O, \theta)$, since $\mathrm{obs}$ does not read the provenance pointer at all --- it follows that $\mathrm{obs}(H_1) = \psi(\{(O,\theta)\}^{\times n}) = \mathrm{obs}(H_2)$.
\end{proof}

Proposition~\ref{prop:presentation-collision} is a direct instance of Theorem~4.1 of \emph{Admissible Reconstruction}: $H_1 \neq H_2$ with $\mathrm{obs}(H_1) = \mathrm{obs}(H_2)$ is precisely the elementary non-identifiability condition, applied here to an observation map defined not for historical linguistics but for the ordinary operation of a memory system deciding what to retain about repeated presentations. The system in Proposition~\ref{prop:presentation-collision} cannot, from its own records, recover whether $n$ occurrences of $(O,\theta)$ reflect one corroborating fact stated once and revisited, or $n$ independent corroborating facts. This is the formal content of the repeated-presentation problem: a memory architecture that defines salience or accessibility as a function of how often $(O,\theta)$ has been recorded, without separately recording $\varepsilon$, is applying Theorem~4.1's collision to its own operation whether or not it is designed with that theorem in mind.

\section{Diagnosis}

The question raised by Proposition~\ref{prop:presentation-collision} is precisely the one Corollary~4.1 of \emph{Admissible Reconstruction} was built to answer: is the resulting ambiguity a correctable failure of this particular observation map, or a structural limit on what any observation map could recover? The corollary's criterion is whether $\reach{H_1} = \reach{H_2}$.

If the admissible continuations available to an observer of the history genuinely do not depend on whether a fact was independently corroborated $n$ times or stated once and referenced $n$ times --- if, for instance, the only thing that matters downstream is whether $(O,\theta)$ was ever asserted at all --- then $\reach{H_1} = \reach{H_2}$, the ambiguity is a structural limit under Theorem~4.2, and no redesign of the observation map that continues to factor through reachable-set structure defined this coarsely can recover the distinction; discarding $\varepsilon$ was not a design flaw but a consequence of what the reachable-set structure was ever going to track. If, on the other hand, independent corroboration and mere repetition license different admissible continuations downstream --- if, for instance, a claim independently corroborated by ten sources warrants stronger downstream commitments than the same claim repeated ten times from a single source, which is the ordinary evidential intuition the repeated-presentation problem trades on --- then $\reach{H_1} \neq \reach{H_2}$, and Corollary~4.1 classifies the failure as correctable: the observation map $\mathrm{obs}$ defined above discarded information, namely $\varepsilon$, that a less lossy observation map could have retained, and a memory architecture that tracks $(O, \theta, \varepsilon)$ jointly rather than $(O,\theta)$ alone would not face this particular collision.

This gives a definite verdict rather than a general injunction to be more careful. The repeated-presentation problem is not, in general, a fundamental limit on machine memory; it is a specific, correctable consequence of choosing an observation map that discards provenance, and Theorem~4.2's own diagnostic --- checking whether provenance affects the reachable set --- is what determines, for any particular downstream use of the memory, whether that discarding was harmless or costly. A memory system is therefore not required to choose once and for all between recording every occurrence at full provenance cost and discarding provenance entirely; Corollary~4.1 supplies the criterion for deciding, use case by use case, which choice the downstream admissibility structure actually requires.
